%Chargement des paquets
\usepackage{amsmath}
\usepackage{amsfonts}
\usepackage{amssymb}
\usepackage{enumerate}
\usepackage{mathtools}
\usepackage{bbm}
\usepackage{xparse, etoolbox}
\usepackage{enumerate}
\usepackage{enumitem}
\usepackage{mathabx}
\usepackage{babel}[french]

%environment
\usepackage{amsthm}
\usepackage{keytheorems}
\usepackage{hyperref} 

%Interface théorème
\newkeytheorem{lem}[
	name = Lemme
]

\newkeytheorem{prop}[
	name = Proposition
]

\newkeytheorem{thm}[
	name = Théorème
]

\newkeytheorem{coro}[
	name = Corollaire
]

\renewcommand*{\proofname}{Démonstration}

\theoremstyle{definition}
\newtheorem{defn}{Définition}

\theoremstyle{definition}
\newtheorem*{nota}{Notation}

\theoremstyle{definition}
\newtheorem*{limi}{Liminaire}

\theoremstyle{remark}
\newtheorem{rmq}{Remarque}

\theoremstyle{plain}
\newtheorem*{fait}{Fait}


%Ensembles de nombres
\newcommand{\N} {\mathbb{N}}
\newcommand{\Ne}{\N^\ast}
\newcommand{\Z} {\mathbb{Z}}
\newcommand{\Q} {\mathbb{Q}}
\newcommand{\R} {\mathbb{R}}
\newcommand{\Rb}{\overline{\mathbb{R}}}
\newcommand{\Rp}{\R_+}
\newcommand{\Rm}{\R_-}
\newcommand{\K} {\mathbb{K}}
\newcommand{\Ket} {\mathbb{K^{\ast}}}
\newcommand{\Cx}{\mathbb{C}}

%Ensembles, tribus
\newcommand{\A}{\mathbb{A}}
\renewcommand{\P}{\mathcal{P}}
\newcommand{\T}{\mathcal{T}}
\newcommand{\F}{\mathcal{F}}
\newcommand{\C} {\mathcal{C}}
\newcommand{\ouv}{\mathcal{O}}
\newcommand{\B}{\mathcal{B}}
\newcommand{\M}{\mathcal{M}}
\newcommand{\etag}{\mathcal{E}}
\newcommand{\etagOp}{\etag^{+}(\Omega)}
\newcommand{\etagO}{\etag(\Omega)}
%%Lp avant quotientage
\newcommand{\Lic}{\mathcal{L}^1}
\newcommand{\Lpc}{\mathcal{L}^p}
\newcommand{\Linfc}{\mathcal{L}^{\infty}}
\newcommand{\Lifull}{\Lic(\Omega, \B, m)}
\newcommand{\Lpfull}{\Lpc(\Omega, \B, m)}
\newcommand{\Linffull}{\Linfc(\Omega, \B, m)}
%%Lp après quotientage
\newcommand{\Li}{L^1}
\newcommand{\Ld}{L^2}
\newcommand{\Lp}{L^p}
\newcommand{\Linf}{L^{\infty}}
\newcommand{\Listd}{\Li(\Omega, m)} 
\newcommand{\Ldstd}{\Ld(\Omega, m)} 
\newcommand{\Lpstd}{\Lp(\Omega, m)} 

%Quelques opérateurs
\DeclareMathOperator{\codim}{codim}
\DeclareMathOperator{\rg}{rg}
\DeclareMathOperator{\tr}{tr}
\DeclareMathOperator{\vect}{Vect}
\DeclareMathOperator{\ima}{Im}
\DeclareMathOperator{\id}{Id}
\DeclareMathOperator{\sign}{sign}
\DeclareMathOperator{\ext}{ext}
\newcommand{\ind}{\mathbbm{1}}
\newcommand*{\spr}[2]{\left(\,#1\,\middle|\,#2\,\right)}
\newcommand{\interior}[1]{%
  {\kern0pt#1}^{\mathrm{o}}%
}
\DeclareMathOperator{\card}{card}
\DeclareMathOperator{\ppcm}{ppcm}

%Commandes de pure flemme
\newcommand{\veps}{\varepsilon}
\newcommand{\intab}{\int_{a}^{b}}
\newcommand{\intR}{\int_{-\infty}^{\infty}}
\newcommand{\intinf}{\int_{- \infty}^{+ \infty}}
\newcommand{\equi}{\Leftrightarrow}
\newcommand{\Kev}{$\K$-espace vectoriel }
\newcommand{\Kevs}{$\K$-espaces vectoriels }
\newcommand{\Rev}{$\R$-espace vectoriel }
\newcommand{\Revs}{$\R$-espaces vectoriels }
\newcommand{\Cev}{$\Cx$-espace vectoriel }
\newcommand{\Cevs}{$\Cx$-espaces vectoriels }
\newcommand{\evn}{(E, \norm{.})}
\newcommand{\interf}[2]{[#1,#2]}
\newcommand{\limb}{\lim\limits_}
\newcommand{\lebn}{\lambda_n}
\newcommand{\pp}{\text{~presque partout}}
\newcommand{\derivp}[2]{\frac{\partial #1}{\partial #2}}
\newcommand{\famiu}{(u_i)_{i \in I}}
\newcommand{\sev}{\text{~s.e.v.}}
\newcommand{\Mnp}{M_{n,p}(\K)}
\newcommand{\Mn}{M_{n}(\K)}
\newcommand{\tq}{\text{~t.q.~}}
\newcommand{\mq}{~montrer~que~}
\newcommand{\et}{\quad \text{et} \quad}
\newcommand{\ou}{\text{~ou~}}
\newcommand{\Kx}{\K[X]}


%Commandes de gars chelou
\renewcommand{\subset}{\subseteq}

%Commande restriction, ty duckduckgo
\def\restr#1#2{\mathchoice
              {\setbox1\hbox{${\displaystyle #1}_{\scriptstyle #2}$}
              \restraux{#1}{#2}}
              {\setbox1\hbox{${\textstyle #1}_{\scriptstyle #2}$}
              \restraux{#1}{#2}}
              {\setbox1\hbox{${\scriptstyle #1}_{\scriptscriptstyle #2}$}
              \restraux{#1}{#2}}
              {\setbox1\hbox{${\scriptscriptstyle #1}_{\scriptscriptstyle #2}$}
              \restraux{#1}{#2}}}
\def\restraux#1#2{{#1\,\smash{\vrule height .8\ht1 depth .85\dp1}}_{\,#2}} 

%Quelques normes
\DeclarePairedDelimiter\abs{\lvert}{\rvert}
\DeclarePairedDelimiter\labs{\bigl\lvert}{\bigr\rvert}
\DeclarePairedDelimiter\norm{\lVert}{\rVert}
\DeclarePairedDelimiterXPP\normi[1]{}\lVert\rVert{_1}{\ifblank{#1}{\:\cdot\:}{#1}}
\DeclarePairedDelimiterXPP\normd[1]{}\lVert\rVert{_2}{\ifblank{#1}{\:\cdot\:}{#1}}
\DeclarePairedDelimiterXPP\normp[1]{}\lVert\rVert{_p}{\ifblank{#1}{\:\cdot\:}{#1}}
\DeclarePairedDelimiterXPP\normq[1]{}\lVert\rVert{_q}{\ifblank{#1}{\:\cdot\:}{#1}}
\DeclarePairedDelimiterXPP\norminf[1]{}\lVert\rVert{_\infty}{\ifblank{#1}{\:\cdot\:}{#1}}

%Bracket en angle
\DeclarePairedDelimiter\angl{\langle}{\rangle}

%Set, parenthèses
\newcommand{\set}[1]{\{ #1 \}}
\newcommand{\bigset}[1]{\bigl\{ #1 \bigr\}}
\newcommand{\Bigset}[1]{\Bigl\{ #1 \Bigr\}}
\newcommand{\bigpar}[1]{\bigl( #1 \bigr)}
\newcommand{\Bigpar}[1]{\Bigl( #1 \Bigr)}

%Opitmisation des opérateurs de comparaison
\DeclarePairedDelimiter\tio{o (}{)}
\DeclarePairedDelimiter\gro{O (}{)}

\newcommand{\Ae}{A^{\ast}}
\newcommand{\Zi}{\Z[i]}
\newcommand{\Zie}{\Z[i]^{\ast}}
\newcommand{\Sd}{\Sigma_2}

\pagestyle{fancy}

\fancyhead[C]{}
\fancyhead[R]{Entiers de Gauss et théorème des deux carrés}

\begin{document}

\underline{Sources :} Rombaldi - Algèbre - page 261 puis 267. \newline

\underline{Un brin d'histoire :} l'anneau des entiers de Gauss a été construit par le mathématicien homonyme afin de résoudre des
problèmes de congruence de carrés, c'est-à-dire afin de résoudre des équations de la forme :
\[ x^2 \equiv y^2 \pmod n , \]
en ignorant les solutions triviales. On verra lors de ce développement que, la structure d'anneau des entiers de Gauss, 
permet de caractériser les entiers naturels qui sont somme de deux carrés.

\vspace{10pt}

On considère comme acquis le résultat suivant :

\begin{thm}[Fermat]
\label{thm:Fermat}
Soit $p$ un entier naturel premier. Pour tout entier relatif $a$ premier avec $p$, on a 
\[ a^{p-1} \equiv 1 \pmod p \]
 et pour tout entier relatif a, on a
\[ a^{p} \equiv a \pmod p \]
\end{thm}

On utilisera aussi le résulat suivant :

\begin{lem}
\label{lem:kdiff}
Pour $k \in \N$, la $k$ème différence de la suite de termes $1^k, 2^k, 3^k, \dots... $ est égale à $k!$.
\end{lem}

\begin{proof}
En notant $\Delta$ l'opérateur différence, défini pour une fonction $f$ par :
\[ \Delta(f)(n) = f(n+1) - f(n) \]
Pour un polynôme $f$ de degré $k$ on a :
\[ \Delta(f)(n) = a_k[(x+1)^k - x^k] + \dots a_1[x+1 - x] \]
où $a_p$ est le coefficient du monome de degré $p$. Or chaque monôme $(x+1)^p - x^p$ est de degré $p-1$, 
donc application de l'opérateur $\Delta$ diminue le degré $k$ du polynôme de 1 et le coefficient de plus haut degré est $k a_k$.
Donc, pour $k$ applications, $\Delta^k(f)$ est une constante égale à $ak!$.
En particulier, pour $f(x)=x^k$, on déduit le résultat.

\end{proof}

\begin{defn}[Stathme]
On appelle stathme sur un anneau commutatif et intègre $A$ une application $\varphi : \Ae to \N$.
\end{defn}

\begin{defn}[Anneau euclidien]
Un anneau commutatif et intègre $A$ est dit euclidien s'il existe un stathme $\varphi$ tel que, pour tout couple $(a,b)$ d'éléments de $A$
avec $b \neq 0$, il existe un couple $(q,r)$ d'éléments de $A$ tel que :
\[ a = bq + r \et r = 0 \ou \varphi(r) < \varphi(b) . \]
\end{defn}

\begin{rmq}
Cette expression permet de définir une division euclidienne.
\end{rmq}

\begin{thm}
Un anneau euclidien est principal. Précisément, pour tout idéal $I$ de $A$ non réduit à $\set{0}$, 
il existe un élément $a_0 \in I \setminus \set{0}$ tel que :
\[ \varphi(a_0) = \min_{a \in I \setminus \set{0}} \varphi(a) \et I = a_0 A . \]
\end{thm}

\begin{proof}
Si $I$ est réduit à 0, il est l'idéal engendré par 0. Sinon, $I \setminus \set{0}$ est non vide et l'image de cet ensemble par $\varphi$
est donc une partie non vide de $\N$, elle admet donc un minimum, noté $n_0$, et il existe $a_0 \in I \setminus \set{0}$ tel que 
$n_0 = \varphi(a_0)$. Soit $x$ un élément quelconque de $I$, en effectuant la division euclidienne de $x$ par $a_0$, on trouve :
\[ x = a_0 q + r \qquad r=0 \ou \varphi(r) < \varphi(a_0) . \]
Or $r \neq 0$ et $\varphi(r) < \varphi(a_0)$ est impossible (car $r = x - a_0 q \in I \setminus \set {0}$), on a donc montré que $x$ est un 
multiple de $a_0$. Comme par ailleurs, tout multiple de $a_0$ est un élément de $I$, on a bien montré que $I = a_0 A$. 
\end{proof}

\begin{nota}
On note $\Zi = \set{a+ib ; (a,b) \in \Z^2}$ l'ensemble des entiers dits de Gauss.
\end{nota}

\begin{thm}
$\Zi$ est un sous-anneau de $\Cx$ stable par l'opération de conjugaison complexe et il est euclidien pour le stathme :
\[ \varphi : u = a+ib \in \Zie \mapsto \abs{u}^2 = a^2 + b^2. \]
\end{thm}

\begin{proof}
Soient $u = a+ ib, v=c+id$ deux éléments de $\Zi$, on a :
\[ u+v = (a+c)+i(b+d) \in \Zi \et uv = (ac - bd) + i(ad +cb) \in \Zi , \]
par ailleurs $1 = 1 + 0i \in \Zi$, donc $\Zi$ est bien un sous-anneau de $\Cx$. 
De plus, $\bar{u} = a-ib \in \Zi$, donc $\Zi$ est stable par l'opération de conjugaison complexe.
Enfin, $\Zi$ hérite des propriétés de commutativité et d'intégrité de $\Cx$. \\
$\varphi$ applique $\Zi$ dans $\N$, c'est donc un stathme sur cet anneau, il reste à démontrer que celui ci est euclidien.
On suppose que $v \neq 0$ et on définit $z = \dfrac{u}{v} = x + iy$ avec $x, y \in \Q$. En utilisant la partition :
\[ \R = \bigcup_{k \in \Z} \left [ k - \dfrac12 ; k + \dfrac12 \right [ , \]
on peut trouver un unique couple d'entiers relatifs $(r,s)$ tel que :
\[ x \in  \left [ r - \dfrac12 ; r + \dfrac12 \right [ \et y \in \left [ s - \dfrac12 ; s + \dfrac12 \right [ . \]
On note $q = r +is \in \Zi$, et on remarque que :
\[ \abs{z - q}^2 = (x-r)^2 + (y-s)^2 \leq \dfrac14 + \dfrac12 < 1 ; \]
ce qui revient à affirmer que $\abs{u-qv}^2 < \abs{v}^2$. On pose donc $r = u -qv \in \Zi$ et on a :
\[ u = qv + r \qquad r=0 \ou \abs{r}^2 < \abs{v}^2 , \]
ce qui démontre le caractère euclidien de $\Zi$.
\end{proof}

\begin{nota}
On note $\Sd = \set { n \in \N ; \exists (a,b) \in \Z^2, n = a^2+b^2}$, l'ensemble des entiers naturels qui s'écrivent comme somme
de deux carrés.
\end{nota}

\begin{lem}
\label{lem:1}
Si $n \in \Sd$ il est alors congru à 0, 1 ou 2 modulo 4. En particulier, si $n$ est impair il est congru à 1 modulo 4.
\end{lem}

\begin{proof}
Il suffit de remarquer que si $a \in \Z$, alors $a^2 = 4n^2$ ou $a^2 = 4(n^2+n)+1$ 
donc $n=a^2+b^2$ est congru à 0, 1 ou 2 modulo 4.
\end{proof}

\begin{lem}
L'ensemble $\Sd$ est stable pour le produit.
\end{lem}

\begin{proof}
Soit $n, m \in \Sd, a = a^2 + b^2$ et $m = c^2 +d^2$. On peut donc écrire $n = \abs{u}$ et $m = \abs{v}$ avec $u, v$ les entiers
de Gauss définis par $u = a+ib, v = c+id$. Or $nm = \abs{uv}$ et $uv \in \Zi$ donc il existe $x, y \in \Z$ tels que $nm = x^2 + y^2$.
\end{proof}

Cette propriété de stabilité nous permet de simplifier le problème : en effet, si tous les facteurs premiers d'un entier sont des éléments
de $\Sd$, alors cet entier sera lui aussi un élément de $\Sd$. 

\begin{thm}
Un nombre premier $p$ est somme de deux carrés si, et seulement si, il égal à 2 ou congru à 1 modulo 4.
\end{thm}

\begin{proof}
Il est clair que $2 \in \Sd$. Si $p>2$, et $p \in \Sd$, alors selon le lemme $\ref{lem:1}$, $p$ est congru à 1 modulo 4. 
Réciproquement, si $p$ est un nombre premier impair congru à 1 modulo 4, alors il existe $n \in \N \tq p=1+4n$
et, selon le petit théorème de Fermat (cf. $\ref{thm:Fermat}$), on peut écrire :
\begin{align*}
1 \equiv 1 \pmod p \\
2^{4n} \equiv 1 \pmod p \\
\dots \\
(4n)^{4n} \equi 1 \pmod p
\end{align*} 
Autrement dit, $p$ divise tous les nombres suivants :
\begin{align*}
2^{4n}-1 & = (2^{2n} + 1)(2^{2n}-1) \\
3^{4n}-2^{4n} & = (3^{2n} + 2^{2n})(3^{2n}-2^{2n}) \\
\dots \\
(4n)^{4n}-(4n-1)^{4n} & = ((4n)^{2n} + (4n-1)^{2n})((4n)^{2n}-(4n-1)^{2n}) .
\end{align*}
Comme $p$ est premier, il divise pour chacun de ces nombres de la forme $AB$, un des deux facteurs $A$ ou $B$.
Supposons que $p$ ne divise que les facteurs $B$, c'est-à-dire $p$ divise :
\[ 2^{2n}-1, \et 3^{2n}-2^{2n}, \quad \dots \et (4n)^{2n}-(4n-1)^{2n} , \]
Alors $p$ divise toutes les différences successives de la suite de nombres $1, 2^{2n}, 3^{2n}, \dots, (4n)^{2n}$,
donc, selon le lemme $\ref{lem:kdiff}$ il divise la $2n$-ème différene successive qui est égale à $(2n)!$. 
Comme $p$ est premier et $p=4n-1$, ceci est absurde.
Donc $p$ divise au moins un facteur $A$, c'est-à-dire qu'il existe $r = 1, \dots, 4n$ tel que $p$ divise $(r^n)^2 + ((r+1)^n)^2$.
Clairement, $a=r^n$ et $b=(r+1)^n$ sont premiers entre eux, 
on peut donc affirmer que $p$ divise une somme de deux carrés premiers entre eux. \\

Supposons maintenant que $p$ soit premier dans $\Zi$, comme cet anneau est principal, et que $p$ divise 
$a^2 + b^2 = (a+ib)(a-ib)$, alors $p$ va diviser l'un ou l'autre de ces facteurs dans $\Zi$, donc va diviser $a$ et $b$ dans $\Z$
ce qui contredit le fait que $a$ et $b$ sont premiers entre eux.
Ainsi, $p$ est réductible dans $\Zi$ et il existe $u, v \in \Zi$ non inversibles et tels que $p=uv$. On a donc :
\[ p^2 = \abs{p}^2 = \abs{u}^2 \abs{v}^2, \]
avec $\abs{u}^2, \abs{v}^2 \in \N \setminus \set{0,1}$, donc $\abs{u}^2=\abs{v}^2=p$, et $p$ est donc somme de deux carrés.
\end{proof}

\begin{lem}
\label{lem:exppair}
Si $n \in \Sd \setminus \set {0}$ admet un diviseur premier $p$ congru à 3 modulo 4, alors $p^2$ divise $n$ et $\dfrac{n}{p^2} \in \Sd$.
\end{lem}

\begin{proof}
Prenons $n$ et $p$ tels qu'annoncés. \\
On a donc $n = a^2 + b^2 \equiv 0 \pmod p$ soit $\bar{a}^2 = - \bar{b}^2$ dans $\Z_p$.
Si on suppose une de ces classes d'équivalence non nulle, on a alors $x = \dfrac{\bar{a}}{\bar{b}}$ (ou $x = \dfrac{\bar{b}}{\bar{a}}$)
qui est solution de l'équation $x^2 = - \bar{1}$ dans $\Z_p$. 
Il s'ensuit que :
\[ (x^2)^{(p-1)/2} \equiv (-1)^{(p-1)/2} \pmod p \et x^{p-1} \equiv 1 \pmod p , \]
Donc $(p-1)/2$ est pair et $p-1 \equiv 0 \pmod 4$.
Comme on a supposé $p$ congru à 3 modulo 4, on a nécessairement $\bar{a}=\bar{b}=0$, donc $p^2$ divise $a$ et $b$ soit :
\[ a = \alpha p \quad ; \quad b = \beta p \et n = p^2(\alpha^2+\beta^2) . \]
\end{proof}

\begin{thm}
Un entier $n \in \Ne$ est somme de deux carrés si, et seulement si, ses éventuels diviseurs premiers congrus à 3 modulo 4 qui
apparaissent dans sa décomposition en facteurs premiers y figurent avec un exposant pair.
\end{thm}

\begin{proof}
Pour la condition nécessaire, on procède par récurrence sur $n \geq 1$. La proposition est trivialement vraie pour $n=1$.
On la suppose vraie pour les entiers de $\Sd \setminus \set{0}$ inférieurs ou égaux à $n-1$, avec $n \in \Sd$ et montrons que la
propriété reste vraie our $n$. Selon le lemme $\ref{lem:exppair}$, si $n$ admet un diviseur premier $p$ congru à 3 modulo 4,
alors $p^2$ apparait dans la décomposition en facteurs premiers de $n$ 
et $\dfrac{n}{p^2} < n$ est un élément de $\Sd \setminus \set{0}$, donc on peut lui appliquer l'hypothèse de récurrence.
On a donc montré que la proposition est héréditaire. \\

La condition suffisante est triviale si on considère la stabilité de $\Sd$ par la multiplication et le fait qu'un nombre premier
congru à 3 modulo 4 aura un carré congru à 1 modulo 4.
\end{proof}

\end{document}